% \begin{savequote}[8cm]
% \textlatin{Neque porro quisquam est qui dolorem ipsum quia dolor sit amet, consectetur, adipisci velit...}

% There is no one who loves pain itself, who seeks after it and wants to have it, simply because it is pain...
%   \qauthor{--- Cicero's \textit{de Finibus Bonorum et Malorum}}
% \end{savequote}

\chapter{\label{ch:6-simplexmaps}Simplex Maps for Multi-Objective Monte Carlo Tree Search} 

    \minitoc













    New Ch5 plan:
    \begin{itemize}
        \item Intro, introduce and motivate simplex maps, motivate by move to weight space, more efficient lookups and ability to use entropy again
        \item Define simplex map interface and implementation
        \item Discussion around the complexity of simplex maps
        \item Use a hypercube example to illustrate that the neighbourhood can still be with respect to dimension and this is seems unnavoidable in general
        \item However, lookups are much more efficient in practise
        \item And sharing values between vertices will cut out a lot of computation when there is no improvement
        \item Algos, BTS and DENTS using simplex maps as a baseline
        \item Empirical results
        \item Theory, contextual regret, non-contextual policies garunteed to have a linear contextual regret
        \item Maybe some theory, probably bootstrap of the ch4 + ch5 results
    \end{itemize}


















\section{Introduction}
\label{sec:6-1-intro}

    \todo{list}
    \begin{itemize}
        \item high level overview of simplex maps work
        \item discuss how simplex maps answer the research questions from introduction chapter
        \item staying in multi-objective land now
        \item Motivated by CHMCTS being slow
    \end{itemize}  

\section{Simplex Maps}
\label{sec:6-2-simplexmaps}

    \todo{list}
    \begin{itemize}
        \item Define simplex map interface
        \item Give details on how to efficiently implement the interface with tree structures
        \item (Good diagram is everything here I think)
    \end{itemize}

\section{Simplex Maps in Tree Search}
\label{sec:6-3-mapsintrees}

    \todo{list}
    \begin{itemize}
        \item Come up with better title for section
        \item Use simplex maps interface to create algorithms from the dents work 
        \item Give a high level idea of what $\delta$ parameter is (used in theory section)
    \end{itemize}

\section{Theoretical Results}
\label{sec:6-4-theory}

    \todo{list}
    \begin{itemize}
        \item Convergence can build ontop of DENTS results
        \item Runtime bounds (better than $O(2^D)$ which is what using convex hulls has)
        \item Simplex map has a diameter $\delta$ (i.e. the furthest away a new context could be from a point in the map)
        \item Bounds can then come from that diameter (which is a parameter of the simplex map/algorithm) and DENTS results
    \end{itemize}

\section{Empirical Results}
\label{sec:6-5-results}

    \todo{list}
    \begin{itemize}
        \item Results from MO-Gymnasium
        \item Compare algorithms using expected utility
    \end{itemize}